\section{数据预处理}
\label{sec:data_preprocessing}

\subsection{数据检查}

销售数据日期范围为 2020-04-01 至 2022-03-31，含 7 家门店、12 个商品代码、13 个商品名称和 8 个类别。商品代码 11001020 对应两个海苔商品名称，经人工确认以商品编号为准，名称差异视为录入笔误。天气数据日期范围为 2020-03-01 至 2022-03-30，不能覆盖销售最后一天 2022-03-31 和未来 2022-04-01 至 2022-04-07。

数据质量检查发现：销售明细中完全重复行 16605 条，经人工确认保留；负销量记录 54 条，作为损耗/冲销类负向调整；交易级数量范围为 -24 至 2712，IQR 规则标记 5536 条交易级极端记录。本文不删除原始记录，而是在处理表中保留正向销量、负向调整和异常标记。主要字段和处理口径见附录 A。

\subsection{日销量面板构造}

本文将交易明细聚合为日期--门店--商品日销量面板，并补齐出现过的门店--商品组合在完整日期上的零销量记录。处理后的日销量面板含 59130 行，日期范围为 2020-04-01 至 2022-03-31；合并外部变量后的建模表同样为 59130 行。

日销量字段包括：
\begin{itemize}
  \item \texttt{daily\_sales}：日净销量，包含负向调整；
  \item \texttt{positive\_sales}：正向销售数量，作为主预测目标；
  \item \texttt{negative\_adjustment\_qty}：负向调整数量绝对值；
  \item \texttt{has\_external\_data}：是否成功匹配附件二外部变量。
\end{itemize}

\subsection{外部变量处理}

附件二的“温度”和“温度.1”经人工确认为最高温和最低温，分别记为 \texttt{max\_temperature} 和 \texttt{min\_temperature}；“活动日”确认为门店促销活动日；“工休日”表示周末/节假日等休息日。附件二未提供湿度字段，因此本文不单独讨论湿度。缺少外部变量的销售日期为 2022-03-31，对应合并后 81 行 \texttt{has\_external\_data=0}。未来 7 天没有附件真实天气和活动日观测，本文在问题四中采用历史同期情景，并明确其条件预测性质。
